\documentclass[11pt,a4paper]{article}
\usepackage[T1]{fontenc}
\usepackage[utf8]{inputenc}
\usepackage{lmodern,microtype}
\usepackage[a4paper,margin=25mm]{geometry}
\usepackage{amsmath,amssymb,amsthm,mathtools}
\usepackage{booktabs,longtable,array}
\usepackage{enumitem}
\usepackage{hyperref}
\hypersetup{colorlinks=true,linkcolor=blue,citecolor=blue,urlcolor=blue,
 pdftitle={AIM 641: research submission},
 pdfauthor={Siavash Sadeghi; AI-assisted research note}}
\usepackage{fancyhdr}
\pagestyle{fancy}
\fancyhf{}
\fancyhead[L]{\small AIM 641 / Siavash Sadeghi}
\fancyhead[R]{\small Research note / 4 October 2026}
\fancyfoot[C]{\thepage}
\setlength{\headheight}{14pt}
\setlength{\parskip}{4pt}
\setlength{\parindent}{0pt}
\setcounter{tocdepth}{1}
\allowdisplaybreaks
\newtheorem{theorem}{Theorem}[section]
\newtheorem{lemma}[theorem]{Lemma}
\newtheorem{proposition}[theorem]{Proposition}
\theoremstyle{remark}
\newtheorem{remark}[theorem]{Remark}
\newcommand{\R}{\mathbb R}
\newcommand{\N}{\mathbb N}
\newcommand{\eps}{\varepsilon}
\newcommand{\norm}[1]{\left\lVert #1\right\rVert}
\newcommand{\abs}[1]{\left\lvert #1\right\rvert}
\newcommand{\gen}{\mathfrak g}
\newcommand{\dd}{\,\mathrm d}
\newcommand{\sgn}{\operatorname{sgn}}
\newcommand{\supp}{\operatorname{supp}}
\newcommand{\dist}{\operatorname{dist}}
\newcommand{\sech}{\operatorname{sech}}

\begin{document}
\title{The cubic Fej\'er inequality in every dimension}
\author{Siavash Sadeghi}
\date{4 October 2026}
\maketitle
\begin{abstract}
This is a claimed proof of the cubic Fejer assertion only. The three quartic assertions grouped in entry 641 are outside the result, so this is a partial resolution of the catalogue entry.
\end{abstract}
\textbf{Provenance and review.} Separated and polished from the supplied AI-assisted research note prepared with ChatGPT. Codex reviewed the mathematical argument and accompanying checks. This is a research claim awaiting independent mathematical review; no independent human audit or formal proof verification is claimed.
\medskip
\section{Entry 641: the cubic Fej\'er inequality in every dimension}\label{sec:cubic}

We settle the cubic assertion within entry 641~\cite{repo641}. Bos proved it through dimension 28 and described a dimension-by-dimension degree-elevation test~\cite{bos}. The point here is to carry out the induction symbolically in the dimension, reducing all remaining dimensions to 45 exact coefficient types.

In barycentric coordinates the prescribed node set consists of the vertices
$e_i$, the ordered edge nodes $t e_i+(1-t)e_j$ for $i\ne j$, where
$t=(1+1/\sqrt5)/2$, and the face centroids $(e_i+e_j+e_k)/3$ for $i<j<k$.
There are $\binom{d+3}{3}$ nodes. Here $\ell_x^{(3,d)}$ denotes the cubic
cardinal polynomial equal to one at $x$ and zero at every other node.

\begin{theorem}\label{thm:cubic}
For every nondegenerate real $d$-simplex, $d\ge1$, Bos's cubic node set $F_{3,d}$ satisfies
\[
 \sum_{x\in F_{3,d}}|\ell_x^{(3,d)}(z)|^2\le1
 \qquad(z\in S_d).
\]
\end{theorem}

\subsection{A dimension-independent symmetric deficit}

Let $M=d+1$ be the number of barycentric coordinates. Work homogeneously with nonnegative $x_1,\ldots,x_M$ and write
\[
 S=\sum_i x_i,\qquad p_j=\sum_i x_i^j.
\]
The homogeneous cubic cardinal polynomials are
\begin{align}
 \ell_i&=\frac{x_i}{2}(12x_i^2-12Sx_i+3S^2-p_2),\label{eq:cubicvertex}\\
 \ell_{ij}&=5x_ix_j\left(\frac{3+\sqrt5}{2}x_i+
                 \frac{3-\sqrt5}{2}x_j-S\right),\quad i\ne j,
                 \label{eq:cubicedge}\\
 \ell_{ijk}&=27x_ix_jx_k,\quad i<j<k.\label{eq:cubicface}
\end{align}
Direct substitution at the three node types verifies cardinality. These are the homogenized polynomials from Bos's construction. Their squared sum, denoted $\mathcal K_M$, is
\begin{align}\label{eq:Ksym}
 \mathcal K_M={}&54p_6+78Sp_5+29S^2p_4-\frac{391}{2}p_2p_4
 -18S^3p_3-144Sp_2p_3+50p_3^2\notag\\
 &+\frac94S^4p_2+\frac{47}{2}S^2p_2^2+\frac{487}{4}p_2^3.
\end{align}
For example, the face contribution follows from
\[
 729\sum_{i<j<k}x_i^2x_j^2x_k^2
 =\frac{243}{2}p_2^3-\frac{729}{2}p_2p_4+243p_6.
\]
The vertex and edge contributions are obtained by the same direct expansion. Formula~\eqref{eq:Ksym} is also independently checked in six variables by the accompanying exact-arithmetic script. This suffices in every dimension: both sides are symmetric homogeneous polynomials of degree six, stable under appending zero coordinates. Every degree-six monomial has support in at most six variables, so its coefficient is determined by a six-variable restriction.

Define the homogeneous deficit
\[
 H_M(x)=S^6-\mathcal K_M(x).
\]
The desired inequality is $H_M\ge0$ on the nonnegative orthant, followed by restriction to $S=1$. Appending a zero coordinate leaves the formula unchanged:
\[
 H_M(x_1,\ldots,x_{M-1},0)=H_{M-1}(x_1,\ldots,x_{M-1}).
\]

\subsection{One induction step covers the whole orthant}

By symmetry, permute any nonnegative $x$ so that $x_M=\min_i x_i$. Set
\[
 z=Mx_M,\qquad \mu_i=x_i-x_M\ge0\quad(1\le i<M).
\]
Then
\begin{equation}\label{eq:centroidsub}
 x_i=\mu_i+z/M\ (i<M),\qquad x_M=z/M.
\end{equation}
Thus the centroid-to-facet substitution does not omit any point of the orthant.

Put $q_j=\sum_{i=1}^{M-1}\mu_i^j$. Under~\eqref{eq:centroidsub},
\begin{equation}\label{eq:pq}
 p_j=\sum_{i=1}^j\binom ji q_i(z/M)^{j-i}+\frac{z^j}{M^{j-1}},
 \qquad 1\le j\le6,
\end{equation}
so $S=q_1+z$. Form the degree-eight polynomial
\begin{equation}\label{eq:Q}
 Q_M(\mu,z)=4M^5(q_1+z)^2
 \left[H_M(\mu_1+z/M,\ldots,\mu_{M-1}+z/M,z/M)
       -H_{M-1}(\mu)\right].
\end{equation}
It is divisible by $z$. Because it is symmetric in the $\mu$ variables, the coefficient of $z^k\mu^\alpha$ depends only on the partition $\alpha$ of $8-k$ consisting of its positive exponents. There are
\[
 p(7)+p(6)+\cdots+p(0)=15+11+7+5+3+2+1+1=45
\]
possible types. For $M\ge8$, every such partition fits within the available $M-1$ variables.

The exact coefficient certificate in Appendix~\ref{app:certificate} establishes that all 45 coefficients are positive for every real $M\ge8$. In particular,
\[
 Q_M(\mu,z)\ge0\qquad(M\ge8,\ \mu,z\ge0).
\]
Except at the origin, the multiplier $4M^5(q_1+z)^2$ is positive. Therefore
\begin{equation}\label{eq:inductionH}
 H_M(\mu+z/M,z/M)\ge H_{M-1}(\mu).
\end{equation}
At the origin both sides vanish.

The cases $1\le d\le6$ are already among Bos's proved cases. The accompanying script checks them again exactly, including the one- and two-dimensional bases and degree elevations $8,5,3,3$ for dimensions $3,4,5,6$, respectively. Starting at $M=7$ and applying~\eqref{eq:inductionH} for $M=8,9,\ldots$ proves Theorem~\ref{thm:cubic}.

\subsection{What the exact verifier checks}

The file \texttt{verify\_cubic.py} uses only exact SymPy expressions and integers or rational numbers in checks that remain active with Python optimization enabled. It checks the direct cardinal-polynomial squared-sum identity; the one-dimensional identity
\[
 H_2(u,v)=12uv(u^2-3uv+v^2)^2;
\]
the two-dimensional positive-coefficient base decomposition; all coefficient types in the four remaining low-dimensional induction steps; and all 45 coefficient polynomials in symbolic $M$ after $M=u+8$.

Its output is:
\begin{quote}\ttfamily\small
PASS: direct cardinal sum-of-squares identity (six variables).\\
PASS: base cases, dimensions one and two.\\
PASS: dimension 3, elevation 8, 123 coefficient types.\\
PASS: dimension 4, elevation 5, 94 coefficient types.\\
PASS: dimension 5, elevation 3, 60 coefficient types.\\
PASS: dimension 6, elevation 3, 64 coefficient types.\\
PASS: all 45 symbolic coefficient types are nonnegative for M >= 8.\\
ALL EXACT CHECKS PASSED.
\end{quote}



\appendix
\section{Complete cubic coefficient certificate}\label{app:certificate}

For a partition $\alpha\vdash8-k$, let $C_{k,\alpha}(M)$ be the coefficient of $z^k\mu_1^{\alpha_1}\cdots\mu_r^{\alpha_r}$ in~\eqref{eq:Q}. The table gives a positive factor $F_{k,\alpha}(M)$ and an array $(a_0,\ldots,a_s)$ such that the following is an exact polynomial identity:
\begin{equation}\label{eq:tablemeaning}
 C_{k,\alpha}(8+u)=F_{k,\alpha}(8+u)\sum_{j=0}^s a_j u^j.
\end{equation}
Every displayed $a_j$ is positive, and each displayed $F(M)$ is positive for $M\ge8$. The empty partition is denoted $\varnothing$; exponents on parts denote repetitions.

\begingroup
\setlength{\tabcolsep}{4pt}
\renewcommand{\arraystretch}{1.13}
\footnotesize
\begin{longtable}{@{}rlll@{}}
\toprule
$k$ & $\alpha$ & $F_{k,\alpha}(M)$ & $(a_0,a_1,\ldots)$\\
\midrule
\endfirsthead
\toprule
$k$ & $\alpha$ & $F_{k,\alpha}(M)$ & $(a_0,a_1,\ldots)$\\
\midrule
\endhead
\bottomrule
\endfoot
1 & $(7)$ & $48 M^{4} \left(M - 1\right)$ & $(1)$\\
1 & $(6,1)$ & $120 M^{4}$ & $(4,1)$\\
1 & $(5,2)$ & $32 M^{4}$ & $(108,19)$\\
1 & $(5,1^{2})$ & $8 M^{4}$ & $(691,111)$\\
1 & $(4,3)$ & $120 M^{4}$ & $(86,13)$\\
1 & $(4,2,1)$ & $8 M^{4}$ & $(1852,327)$\\
1 & $(4,1^{3})$ & $24 M^{4}$ & $(1681,231)$\\
1 & $(3^{2},1)$ & $528 M^{4}$ & $(37,7)$\\
1 & $(3,2^{2})$ & $16 M^{4}$ & $(246,191)$\\
1 & $(3,2,1^{2})$ & $16 M^{4}$ & $(2933,543)$\\
1 & $(3,1^{4})$ & $48 M^{4}$ & $(3133,438)$\\
1 & $(2^{3},1)$ & $144 M^{4}$ & $(144,59)$\\
1 & $(2^{2},1^{3})$ & $48 M^{4}$ & $(2773,493)$\\
1 & $(2,1^{5})$ & $240 M^{4}$ & $(1603,234)$\\
1 & $(1^{7})$ & $30240 M^{4}$ & $(29,4)$\\
2 & $(6)$ & $36 M^{3} \left(M - 1\right)$ & $(22,3)$\\
2 & $(5,1)$ & $12 M^{3}$ & $(1886,593,43)$\\
2 & $(4,2)$ & $4 M^{3}$ & $(23766,6679,461)$\\
2 & $(4,1^{2})$ & $4 M^{3}$ & $(46558,13169,879)$\\
2 & $(3^{2})$ & $8 M^{3}$ & $(19494,5287,359)$\\
2 & $(3,2,1)$ & $16 M^{3}$ & $(14394,4748,355)$\\
2 & $(3,1^{3})$ & $24 M^{3}$ & $(29968,8563,576)$\\
2 & $(2^{3})$ & $24 M^{3}$ & $(34,1660,193)$\\
2 & $(2^{2},1^{2})$ & $8 M^{3}$ & $(65834,23543,1817)$\\
2 & $(2,1^{4})$ & $72 M^{3}$ & $(24308,7179,498)$\\
2 & $(1^{6})$ & $720 M^{3}$ & $(5816,1611,108)$\\
3 & $(5)$ & $4 M^{2} \left(M - 1\right)$ & $(2538,711,47)$\\
3 & $(4,1)$ & $4 M^{2}$ & $(117366,49147,6686,295)$\\
3 & $(3,2)$ & $8 M^{2}$ & $(102258,42029,5739,258)$\\
3 & $(3,1^{2})$ & $8 M^{2}$ & $(234260,99829,13674,605)$\\
3 & $(2^{2},1)$ & $8 M^{2}$ & $(159730,81927,12648,611)$\\
3 & $(2,1^{3})$ & $24 M^{2}$ & $(177980,80093,11338,515)$\\
3 & $(1^{5})$ & $480 M^{2}$ & $(21806,9274,1263,56)$\\
4 & $(4)$ & $M \left(M - 1\right)$ & $(130056,52702,6921,295)$\\
4 & $(3,1)$ & $2 M$ & $(1844808,1015586,206723,18410,605)$\\
4 & $(2^{2})$ & $2 M$ & $(1571928,903466,192249,17862,611)$\\
4 & $(2,1^{2})$ & $2 M$ & $(4011400,2362530,501963,45998,1545)$\\
4 & $(1^{4})$ & $24 M$ & $(792080,454696,94306,8465,280)$\\
5 & $(3)$ & $2 \left(M - 1\right)$ & $(420984,224198,44113,3800,121)$\\
5 & $(2,1)$ & $6$ & $(2183112,1536018,425265,57927,3887,103)$\\
5 & $(1^{3})$ & $12$ & $(2386440,1698730,470883,63852,4255,112)$\\
6 & $(2)$ & $M - 1$ & $(325512,179338,36255,3184,103)$\\
6 & $(1^{2})$ & $2$ & $(2335608,1664406,463769,63260,4237,112)$\\
7 & $(1)$ & $8 \left(M - 3\right) \left(M - 1\right)$ & $(2376,874,103,4)$\\
8 & $\varnothing$ & $\left(M - 3\right) \left(M - 1\right)$ & $(2376,874,103,4)$\\
\end{longtable}
\endgroup

Here is an exact way of checking the table without expanding in an unbounded number of variables. Treat $q_1,\ldots,q_6,z,M$ as formal indeterminates and substitute~\eqref{eq:pq} into~\eqref{eq:Q}. For a product $q_{b_1}\cdots q_{b_s}$, define
\[
 T((b_1,\ldots,b_s),\alpha)
 =[\mu^\alpha]\prod_{j=1}^s\left(\sum_i\mu_i^{b_j}\right).
\]
It is computed by the finite recursion
\begin{align*}
 T((),\alpha)&=\mathbf1_{\alpha=0},\\
 T((b_1,\ldots,b_s),\alpha)
 &=\sum_{i:\alpha_i\ge b_1}
 T((b_2,\ldots,b_s),\alpha-b_1e_i).
\end{align*}
This counts all assignments of the power-sum factors to the target monomial's variables, including the correct multiplicities for repeated factors. Apply it to each formal monomial in $Q_M$. The resulting coefficients are the 45 polynomials represented in the table. The accompanying JSON file also gives the full coefficients of each $C_{k,\alpha}(8+u)$ without factor removal.

The script \path{check_printed_table.py} also compares the printed table symbolically with the generated certificate. The separate file \path{crosscheck_cubic.py} expands the polynomial directly in a multivariate polynomial ring at $M=8$. All 3,432 nonzero monomial coefficients agree exactly with the 45-type certificate. This cross-check does not use the power-sum coefficient recursion.

The all-dimensional argument is a uniform symbolic certificate, not a check at a finite list of large dimensions. The quartic conjectures in the same catalogue entry are not part of Theorem~\ref{thm:cubic}.

\begin{thebibliography}{9}
\bibitem{repo641}
M. Colbrook, AIM catalogue, \href{https://github.com/MColbrook/AIM/blob/59a8f0c2957dd272f56bcf282dfe0c066a15f441/problems/641-bos-simplex-interpolation-nodes.md}{entry 641: Bos's cubic and quartic simplex interpolation nodes}. Repository statement checked 4 October 2026.

\bibitem{bos}
L. Bos, \emph{On Fekete points for a real simplex}, Indagationes Mathematicae \textbf{34} (2023), 274--293. \href{https://arxiv.org/abs/2205.06498}{arXiv:2205.06498}; \href{https://doi.org/10.1016/j.indag.2022.11.003}{doi:10.1016/j.indag.2022.11.003}. Section 3 and Proposition 3.1.
\end{thebibliography}
\end{document}
